%% =============================================================================
%% sm_d2_informativeness.tex — SM-D, Subsection D.2: Informativeness Assessment
%% Body content only (\subsection level); parent structure in sm_d.tex
%% Source: dev/stan/R/81_informativeness_tests.R
%% Corrections applied: T3-1 (full-covariate informativeness tests),
%%   T3-2 (Hausman test with survey::svyglm + cluster-robust SEs)
%% =============================================================================

\subsection{Informativeness assessment}
\label{smd:informativeness}

A sampling design is \emph{informative} if the distribution of the
outcome among sampled units differs from the superpopulation
distribution, even after conditioning on covariates
\citep{Pfeffermann1993}.  Under informative sampling, unweighted
likelihood-based inference targets a sample-specific, not population,
quantity, and pseudo-likelihood weighting becomes essential for
consistent estimation~\citep{SavitskyToth2016}.  We assess
informativeness of the NSECE~2019 design through three complementary
diagnostics.

{\upshape\textbf{(Correction T3-1.)}}
The original informativeness tests conditioned on a single covariate
(poverty rate).  Since the full model includes $P = 5$ covariates per
margin, omitted variable bias may mask or inflate apparent
informativeness.  We now report both the original univariate and the
corrected full-covariate specifications.

{\upshape\textbf{(Correction T3-2.)}}
The original Hausman comparison used \texttt{lm(weights = w)} for the
weighted estimator, which provides only heteroscedasticity-robust
standard errors.  We now use \texttt{survey::svyglm()} with the full
stratified cluster design, yielding cluster-robust standard errors
that account for within-PSU correlation.

\paragraph{Test~1: Marginal weight--outcome correlations.}
The correlation between $z_i$ (serves infants) and the survey weight
$w_i$ is $\corr(z_i, w_i) = -0.041$; the correlation between
the IT share $Y_i / n_i$ (among servers) and $w_i$ is
$\corr(Y_i/n_i, w_i) = -0.003$.  Both are negligibly small,
suggesting limited \emph{unconditional} informativeness.  However,
marginal correlations can mask conditional dependence when the design
oversamples on covariates correlated with the outcome, motivating the
regression-based test below.

\paragraph{Test~2: Regression-based informativeness test.}
Following~\citet{Pfeffermann1993}, we test $H_0\colon \gamma = 0$ in
the augmented model
\begin{equation}\label{smd:eq-pfeffermann}
  g(y_i) = \bx_i^\t \bbeta + \gamma \log(w_i) + \varepsilon_i,
\end{equation}
where $g(\cdot)$ is the logit link for the extensive margin and the
identity for the intensive margin.  If the design is non-informative,
$\log(w_i)$ carries no additional predictive power for $y_i$ beyond
$\bx_i$, so $\gamma = 0$.

\Cref{smd:tab-pfeffermann} reports the estimated $\hat{\gamma}$ and
its $p$-value under two covariate specifications: (i)~univariate, with
only the standardized poverty rate; and (ii)~full, with all four
community-level predictors (poverty, urbanicity, percent Black, percent
Hispanic).

\begin{table}[htbp]
  \centering
  \caption{Regression-based informativeness test~\citep{Pfeffermann1993}.
    The augmented model \eqref{smd:eq-pfeffermann} adds $\log(w_i)$ to the
    covariate vector; $\hat{\gamma}$ is the coefficient on $\log(w_i)$
    with its two-sided $p$-value.  {\upshape\textbf{(Correction T3-1:)}}
    the full-covariate specification reverses the extensive-margin
    finding from non-significant to significant.}
  \label{smd:tab-pfeffermann}
  \smallskip
  \begin{tabular}{llrr}
    \toprule
    Specification & Margin & $\hat{\gamma}$ & $p$-value \\
    \midrule
    Univariate (poverty only)
      & Extensive  & $0.020$  & $0.369$   \\
      & Intensive  & $0.009$  & $0.002$   \\[4pt]
    Full covariates (T3-1)
      & Extensive  & $0.086$  & $<0.001$  \\
      & Intensive  & $0.006$  & $0.050$   \\
    \bottomrule
  \end{tabular}
\end{table}

Two features of~\cref{smd:tab-pfeffermann} merit comment.  First, under
the univariate specification, the extensive margin appears
non-informative ($p = 0.369$) while the intensive margin is clearly
informative ($p = 0.002$).  Second, adding the remaining three
covariates \emph{reverses} the extensive-margin verdict: $\hat{\gamma}$
increases fourfold (from $0.020$ to $0.086$) and becomes highly
significant ($p < 0.001$).  This reversal arises because the univariate
model confounds the direct weight--outcome relationship with indirect
pathways through the omitted covariates (urbanicity, racial
composition).  When these confounders are included, the residual
association between $\log(w_i)$ and participation becomes detectable.
The intensive margin remains informative under both specifications,
though the coefficient attenuates slightly (from $0.009$ to $0.006$)
as the full model absorbs some of the weight--outcome association.

The finding that both margins are informative under the full
specification reinforces the case for pseudo-likelihood weighting in the
Bayesian hierarchical model and for the sandwich variance correction
developed in~\cref{app:survey}.

\paragraph{Test~3: Hausman-type comparison.}
As a complementary diagnostic, we compare unweighted and weighted
coefficient estimates for each margin.  Under non-informative sampling,
both estimators are consistent for the same population parameters, but
the weighted estimator is less efficient; a large discrepancy signals
that the design distorts the effective covariate--outcome relationship.

{\upshape\textbf{(Correction T3-2:)}} the weighted estimator uses
\texttt{survey::svyglm()} with the full stratified cluster design ($H =
30$ strata, $C = 415$ PSUs), which produces cluster-robust standard
errors.  The Hausman test statistic for each parameter is
\begin{equation}\label{smd:eq-hausman}
  Z_k = \frac{\hat{\beta}_k^{\mathrm{WT}} -
    \hat{\beta}_k^{\mathrm{UW}}}
  {\sqrt{\widehat{\Var}(\hat{\beta}_k^{\mathrm{WT}}) -
    \widehat{\Var}(\hat{\beta}_k^{\mathrm{UW}})}},
\end{equation}
which is asymptotically standard normal under $H_0$ (no
informativeness); the variance difference in the denominator is valid
because the weighted estimator has weakly larger variance under
non-informativeness.

\begin{table}[htbp]
  \centering
  \caption{Hausman-type comparison of unweighted and
    design-weighted coefficient estimates.
    {\upshape\textbf{(Correction T3-2:)}} the weighted estimator uses
    \texttt{survey::svyglm} with cluster-robust standard errors.
    $Z$ is the per-parameter Hausman statistic
    \eqref{smd:eq-hausman}.  Asterisks: ${*}$ denotes $p < 0.10$,
    ${**}$ denotes $p < 0.05$.}
  \label{smd:tab-hausman}
  \smallskip
  \begin{tabular}{llrrrrr}
    \toprule
    Margin & Parameter
      & $\hat{\beta}^{\mathrm{UW}}$ & $\hat{\beta}^{\mathrm{WT}}$
      & $Z$ & $p$-value & \\
    \midrule
    \multirow{5}{*}{Extensive}
      & Intercept & $0.614$  & $0.662$  & $0.82$  & $0.41$ & \\
      & Poverty   & $-0.158$ & $-0.199$ & $-0.64$ & $0.52$ & \\
      & Urban     & $0.176$  & $0.234$  & $2.02$  & $0.044$ & $^{**}$ \\
      & Black     & $-0.044$ & $0.086$  & $1.89$  & $0.059$ & $^{*}$ \\
      & Hispanic  & $-0.043$ & $0.059$  & $1.51$  & $0.13$ & \\[4pt]
    Intensive
      & \multicolumn{5}{l}{\emph{all $p > 0.24$; none significant}} & \\
    \bottomrule
  \end{tabular}
\end{table}

\Cref{smd:tab-hausman} reveals a nuanced pattern.  On the extensive
margin, the poverty coefficient itself does not differ significantly
between estimators ($Z = -0.64$, $p = 0.52$), but the
\emph{urbanicity} coefficient shows a significant shift
($\hat{\beta}^{\mathrm{UW}} = 0.176$ versus
$\hat{\beta}^{\mathrm{WT}} = 0.234$; $Z = 2.02$, $p = 0.044$), and
the Black composition coefficient is marginally significant ($Z = 1.89$,
$p = 0.059$).  This indicates that the NSECE oversampling of
high-need areas primarily distorts the estimated effect of urbanicity
on IT participation---consistent with the fact that the sampling frame
explicitly targets urban and high-poverty PSUs.  On the intensive
margin, no individual parameter shows a significant Hausman
discrepancy (all $p > 0.24$), suggesting that the informativeness
detected by Test~2 is diffuse rather than concentrated on a single
covariate.

\paragraph{Summary.}
The three diagnostics paint a consistent picture.  Marginal correlations
(Test~1) are negligible, but the regression-based test (Test~2) reveals
statistically significant informativeness on \emph{both} margins once
the full covariate vector is conditioned upon---a finding that was
masked in the original univariate specification.  The Hausman comparison
(Test~3) localises the extensive-margin informativeness to the
urbanicity and racial composition covariates rather than poverty itself.
Together, these results justify two design choices in the main
analysis (\cref{sec:application}): (i)~incorporating survey weights
via pseudo-likelihood, and (ii)~applying the sandwich variance
correction (\cref{sec:survey}) and the Cholesky affine calibration
(\cref{app:survey}) to ensure that posterior credible intervals
reflect the true design-adjusted uncertainty.
